\section{Introduction}
% Timely and formative feedback is crucial in academia. It allows students to reinforce learning, increase their knowledge, and develop essential skills for more advanced topics. However, manual assessment of student code submissions is time-intensive and delays feedback~\cite{barczak2023automated,10.1145/3591210}. Such delays negatively impact student learning and performance, particularly in subsequent assignments.

% Instructors often address this issue by solely focusing on code correctness and functionality while neglecting other coding aspects, such as code quality~\cite{Chen2018AnAA,9814913}. However, overlooking such crucial coding aspects may encourage students to lean into writing code that functions but with poor code quality, increasing the risk of bugs and maintenance costs%~\cite{10.1145/1095714.1095770}
% . Even worse, instructors may address it by employing multiple teaching assistants in high-enrollment courses, leading to inconsistent grading and feedback~\cite{10.1145/3591210}. Therefore, these challenges require a more efficient and standardized approach to assessing programming assignments.

% In this paper, we introduce~\textbf{\codeinspectoren}---an automated code assessment system for CS1-level Java programming assignments. \codeinspector is a scalable and modular solution that leverages open-source tools such as GitHub, Jenkins~\cite{Jenkins}, Apache Ant~\cite{ant}, JUnit~\cite{junit}, CheckStyle~\cite{checkstyle}, PMD~\cite{pmd}, and Large Language Models (LLMs) to streamline key aspects of the assessment process, including code quality analysis and unit test generation and execution. It provides students with immediate, formative feedback on correctness, functionality, and code style while reducing instructors' assessment workload.


% The contributions of this paper include the following: 
% \begin{enumerate}
%     \item A scalable and modular automated assessment system leverages open-source tools, allowing customization for various coding aspects and assessment needs.
%     \item The integration of an LLM for automated test case generation, enhancing the efficiency of assessment workflows.
%     \item A grading scheme that ensures fair assessment by weighting penalties based on error density and severity rather than the total number of violations.
% \end{enumerate}

% { The remainder of this paper is organized as follows. 
% Section~\ref{sec:background} presents related work relevant to our research. 
% Section~\ref{sec:sysdesign} describes the proposed system design.
% Section~\ref{sec:evaluation} presents the evaluation results and analysis.
% Finally, Section~\ref{sec:conclusions} concludes the paper and outlines directions for future work.
% }